\begin{figure}[htbp]
\centering
\begin{tikzpicture}[>=Latex,font=\small]
 \fill[blue!5] (3,0) rectangle (6,2.05);
 \draw[black!55] (0,0) rectangle (3,2.05);
 \draw[black!55] (3,0) rectangle (6,2.05);
 \draw[black!55] (6,0) rectangle (9,2.05);
 \draw[teal!80!black,line width=2pt] (3,0)--(3,2.05);
 \node[align=center,text width=2.6cm] at (1.5,1.55) {upstream\\neighbour $N$};
 \node[align=center,text width=2.6cm] at (4.5,1.55) {inspected\\cell $K$};
 \node[align=center,text width=2.6cm] at (7.5,1.55) {downstream\\neighbour};
 \node at (1.5,.6) {$u_N$};
 \node at (4.5,.6) {$u_K$};
 \draw[->,thick] (4.1,2.65)--(5.25,2.65);
 \node[above] at (4.68,2.65) {transport $\boldsymbol v$};
 \draw[->,thick] (3,1.05)--(2.28,1.05);
 \node[below] at (2.53,1.05) {$\boldsymbol n_K$};
 \foreach \y in {.25,.5,.75,1.3,1.8} {
   \fill[teal!80!black] (3,\y) circle (1.4pt);
 }
 \draw[teal!80!black] (3,-.08)--(3,-.45);
 \node[below,align=center,text=teal!80!black] at (3,-.45)
   {inflow: $\boldsymbol v\cdot\boldsymbol n_K<0$\\integrate $u_K-u_N$ here};
 \draw[black!55] (6,-.08)--(6.7,-.45);
 \node[below,align=center] at (7,-.45)
   {outflow for $K$\\not in its inflow integral};
\end{tikzpicture}
\caption{KXRCF on an interior patch with uniform rightward transport.
For the inspected cell $K$, the left face is inflow, the right face is
outflow, and the horizontal faces are tangential. The highlighted face enters
$K$'s signed jump integral. Colours identify the inspected cell and face,
not a computed troubled-cell classification. For curved or varying flow,
the sign is tested separately at each face quadrature point.}
\label{fig:kxrcf-inflow}
\end{figure}
